\section{Related Work}

% \subsubsection{GPU Graph Processing beyond Device Memory}
% TODO: Add the original Medusa, CuSha, and Gunrock papers/bibliography entries. % Medusa和Cusha没添加，觉得参考文献有点时间旧
% The massive parallelism of GPUs makes them a natural platform for vertex-centric graph processing. Early frameworks map this computation to GPU threads and organize parallel traversal and value propagation~\cite{2016-Gunrock} [Medusa, CuSha, Gunrock]. However, their scalability is constrained by device memory capacity, as large graphs can exceed the memory available on a single GPU and require out-of-memory (OOM) processing.
% For static graphs, OOM systems retain graph data in host memory and transfer the portions needed for GPU computation. Subway constructs and transfers active subgraphs, EMOGI enables fine-grained zero-copy reads of host-resident edges, and HyTGraph selects between explicit transfers and zero-copy access~\cite{2020-Subway, 2020-EMOGI, 2023-HyTGraph} [Subway, EMOGI, HyTGraph]. CGgraph further combines a reusable GPU-resident subgraph with CPU-GPU task allocation for cooperative execution~\cite{2024-CGgraph} [CGgraph]. These systems avoid loading the entire graph onto the GPU; however, since PCIe transfer speeds are significantly lower than memory access speeds, reading data from the host still incurs overhead.
% The dynamic graph processing system introduces the requirement of maintaining adjacency storage as updates arrive. GPU structures such as Hornet and faimGraph support topology changes within device memory~\cite{2018-Hornet, 2018-faimGraph} [Hornet, faimGraph]. Extending GPU processing to OOM dynamic graphs requires both access to host-resident topology and efficient maintenance of graph data and index views on the GPU after updates. In contrast to static graph storage, the host representation and any GPU-resident adjacency data must remain consistent with the current update batch before computation proceeds. In streaming graphs, even though only a portion of the graph data is updated, recalculating the results after every update would significantly increase communication overhead. Therefore, merely addressing capacity constraints does not eliminate the overhead associated with continuous graph analysis.

\subsection{Dynamic Graph Processing on CPUs}

CPU-based dynamic graph processing has primarily focused on efficient graph updates and incremental analytics. For graph storage, Sortledton, Spruce, DCSR, Bubble, and LSMGraph develop representations that support updates while retaining efficient access for analytics~\cite{2022-Sortledton,2024-Spruce,2025-DCSR,2025-Bubble,2024-LSMGraph}.
RisGraph combines graph updates with low-latency incremental analysis~\cite{2021-RisGraph}. On the computation side, a small update often affects only a subset of the existing results. KickStarter, GraphBolt, and DZiG exploit this property to avoid recomputation over unaffected parts of the graph~\cite{2017-KickStarter, 2019-GraphBolt,2021-DZiG}.
Tripoline, CommonGraph, and IncBoost further extend incremental processing to different query, snapshot, and update
settings~\cite{2021-Tripoline,2023-CommonGraph,2024-IncBoost}. These techniques reduce redundant CPU computation, motivating efforts to exploit GPU parallelism for evolving graphs.

\subsection{Dynamic Graph Processing on GPUs}

GPU graph processing initially focused on static graphs. Gunrock, for example, provides parallel primitives for graph
analytics~\cite{2016-Gunrock}. For static graphs that exceed GPU memory, Subway transfers active subgraphs, EMOGI accesses host-resident edges, and HyTGraph adapts the data access strategy to the workload~\cite{2020-Subway,2020-EMOGI,2023-HyTGraph}. CGgraph retains a reusable subgraph on the GPU and distributes work between the CPU and GPU~\cite{2024-CGgraph}. Processing evolving graphs adds further requirements: the graph representation must support
updates without impeding traversal, while incremental execution must handle irregular and varying amounts of affected work.

GPU-resident dynamic graph systems primarily address the update and traversal tradeoff. GPMA supports parallel batches of edge updates using packed memory arrays~\cite{2017-GPMA}. Hornet and faimGraph develop dynamic GPU graph structures~\cite{2018-Hornet, 2018-faimGraph}, while GASgraph uses a SubHPMA-based representation for streaming graph
processing~\cite{2025-GASgraph}. Their reliance on GPU-resident graph data limits the graph sizes they can accommodate. For graphs exceeding GPU memory, Grapin combines GPU incremental computation with hot-subgraph caching to reduce accesses to host-resident data~\cite{2025-Grapin}. CGCGraph co-executes concurrent snapshot analyses by processing shared graph data on the GPU and unshared data on the CPU~\cite{2025-CGCGraph}. POEGA uses proxy graphs to guide out-of-memory refinement and processes multiple snapshots concurrently~\cite{2026-POEGA}. CoIncGraph instead partitions computation affected by streaming updates between the CPU and GPU, coupling incremental execution with cooperative processing of graphs that exceed GPU memory.

% 这个上面是否要写related work，到时候依据篇幅来定吧，看Background章节写了一些Related work，所以这块暂时不安排也是可以的
\section{Conclusion}
% 暂时整了一个初稿，后面再根据结果等做进一步的优化
We presented CoIncGraph, a CPU--GPU cooperative system for incremental processing of out-of-memory streaming graphs. Its heterogeneous source-isolated graph organization preserves topological locality, limiting adjacency relocation and GPU index maintenance to changed sources. Collaborative deletion repair and insertion propagation restrict computation to affected vertices while coordinating CPU preparation with GPU execution. Ordered propagation and shared batch preparation further reduce redundant edge scans and CPU maintenance overhead. Experiments on six real-world graphs demonstrate that CoIncGraph outperforms Ingress, SHGraph, and Grapin. These results highlight the benefits of jointly exploiting update locality and CPU--GPU cooperation for efficient streaming graph processing.